%% kapa: exam=1 round=21 year=2010 period=2 form=A subject=appraisal_law subject_ko=부동산관계법규 seq=16 number=56
%% source: Q-Net 감정평가사 기출문제 (cst00309) artlSeq=5024181 file=extracted/1008/5024181/2010년 제21회 감정평가사 1차시험 기출문제/2교시.hwp
%% conversion: kordoc 4.24.51 -> tools/gen1.py (auto); review=proofread-gpt-sol; flags=none
\begin{question}{56}
국토의 계획 및 이용에 관한 법령상 토지거래허가에 관한 설명으로 옳지 않은 것은?
\begin{choices}
  \item 토지의 투기적인 거래가 성행할 우려가 있는 지역에 대하여는 토지거래계약에 관한 허가구역으로 지정할 수 있다.
  \item 「공익사업을 위한 토지 등의 취득 및 보상에 관한 법률」에 따라 토지를 협의취득할 때에도 토지거래계약에 관한 허가를 받아야 한다.
  \item 토지거래계약에 관한 허가구역에 있는 토지에 관한 소유권을 이전하는 계약을 체결하려는 당사자는 공동으로 시장․군수 또는 구청장의 허가를 받아야 한다.
  \item 토지거래계약에 관한 허가구역 안에서 토지거래계약허가를 받지 아니하고 체결한 토지거래계약은 그 효력이 없다.
  \item 토지거래계약에 관한 허가신청에 대하여 불허가처분을 받은 자는 그 통지를 받은 날부터 1개월 이내에 해당지역의 시장․군수 또는 구청장에게 해당 토지에 관한 권리의 매수를 청구할 수 있다.
\end{choices}
\end{question}
\answer{2}
